% Compact replacement section for the AE response.
% Primary results: wl_code_package_20261015/results and
% wl_leave_one_type_out_20261015; cumulative WL depth h=2 with effective bounds
% retained as variable restrictions and excluded from the graph.
% The surrounding reply letter should load amsmath, booktabs, float,
% natbib, and xcolor, and define the color replyblue.
\providecommand{\refdata}{\textsc{Ref-Data}}
\providecommand{\testdata}{\textsc{Large-Scale-OR}}
\newcommand{\LK}[1]{\textcolor{blue}{#1$^{\text{LK}}$}}

\section{Structural Similarity and Robustness}
\label{sec:ae-structural-similarity-compact}

\begingroup\color{replyblue}
\textbf{Response.}
We agree that a shared problem category does not imply a shared
formulation. We therefore separate two questions: (i) how close each
formulation is to the available reference pool under an explicit graph
encoding, and (ii) whether the pipeline still performs when examples,
routes, or task requirements are changed.

\medskip
\textbf{Motivation and method.}
We represent each linear formulation by its variable--constraint
bipartite graph, following the MILP graph representation of
\citet{Gasse2019}, and compare labeled graphs with the Weisfeiler--Lehman
(WL) subtree construction of \citet{Shervashidze2011}. This provides an
untrained and auditable structural diagnostic; related work also uses
graph-based MILP similarity for solver configuration
\citep{Hosny2024}. The score is a diagnostic of the selected graph
representation, not a learned measure of business-level similarity.

The current code package recomputes 15 reference formulations, 101
Large-Scale-OR formulations, and 36 task variants from the original LP
files. A singleton row is converted to an effective variable bound and
intersected with native bounds, so equivalent storage in ``Bounds'' and
``Subject To'' does not create an artificial graph difference. In the
primary representation, these effective bounds remain active in the
optimization model but do not create graph nodes. Variable domains, all
remaining rows, and isolated variables are retained; no general
presolve or variable elimination is applied. This normalization is not a
complete canonicalization of mathematically equivalent formulations.

Each variable and retained row is a node, with an edge for every nonzero
matrix coefficient. Initial labels distinguish binary, integer, and
continuous variables ($\texttt{V:B}$, $\texttt{V:I}$, $\texttt{V:C}$)
and equality versus inequality rows ($\texttt{R:E}$, $\texttt{R:I}$).
Coefficient magnitudes and signs, objective terms, right-hand sides,
bound values, and names are omitted. Thus the diagnostic emphasizes
variable--constraint incidence and domains while deliberately omitting
some sources of within-category heterogeneity.

\medskip
\textbf{WL features and reference matching.}
Let $l_0(v)$ be the initial label. At round $i\geq1$, all nodes update
synchronously using
\[
M_i(v)=\{\!\{l_{i-1}(u):u\in\mathcal N(v)\}\!\},\qquad
s_i(v)=\operatorname{concat}\!\left(l_{i-1}(v),
\operatorname{sort}(M_i(v))\right),\qquad
l_i(v)=f(s_i(v)),
\]
where the same dictionary $f$ assigns one label to equal signatures
across all graphs. We concatenate label counts from rounds $0$--$2$ and
use cosine-normalized cumulative counts:
\[
S_2(P,Q)=
\frac{\displaystyle\sum_{i=0}^{2}\sum_{\sigma\in\Sigma_i}
c_i(G_P,\sigma)c_i(G_Q,\sigma)}
{\displaystyle\sqrt{\left(\sum_{i=0}^{2}\sum_{\sigma\in\Sigma_i}
c_i(G_P,\sigma)^2\right)
\left(\sum_{i=0}^{2}\sum_{\sigma\in\Sigma_i}
c_i(G_Q,\sigma)^2\right)}}.
\]
For each test $P$, we report the maximum over its designated
same-category pool, $S_{\rm same}(P)$, and the maximum after removing
all references with the test's category, $S_{\rm LOTO}(P)$. Thus a
Mixture test removes the eight Mixture references but retains the two
Others references; an Others test removes only the two Others
references. Mixture and Others remain separate reference pools when
computing $S_{\rm same}$.

\begin{table}[H]
\centering\scriptsize
\caption{Primary WL results and structural LOTO scores. All scores use
cumulative rounds $0$--$2$, effective bounds excluded from the graph, and
cosine normalization. Each score is a median or mean of per-instance
maxima as indicated.}
\label{tab:ae-wl-compact-results}
\setlength{\tabcolsep}{4.5pt}
\begin{tabular}{@{}llrrrrr@{}}
\toprule
Set & Category & Tests & Median $S_{\rm same}$ & Mean $S_{\rm same}$ &
Median $S_{\rm LOTO}$ & Mean $S_{\rm LOTO}$\\
\midrule
Original & NRM     & 25 & 0.3234 & 0.3234 & 0.3330 & 0.3330\\
Original & RA      & 22 & 0.6562 & 0.6564 & 0.6500 & 0.6505\\
Original & TP      & 9  & 0.6240 & 0.6314 & 0.6327 & 0.6369\\
Original & AP      & 5  & 0.0508 & 0.0551 & 0.4424 & 0.4437\\
Original & UFLP    & 14 & 0.5503 & 0.5557 & 0.4446 & 0.4493\\
Original & Mixture & 18 & 0.3983 & 0.4499 & 0.4057 & 0.4327\\
Original & Others  & 8  & 0.2584 & 0.3025 & 0.3615 & 0.4070\\
\midrule
Original & All     & 101& 0.5065 & 0.4632 & 0.4424 & 0.4745\\
\midrule
Variant & AP      & 3  & 0.3569 & 0.3582 & 0.2749 & 0.2752\\
Variant & UFLP    & 3  & 0.1981 & 0.1972 & 0.3078 & 0.3079\\
Variant & Mixture & 17 & 0.6665 & 0.6157 & 0.2520 & 0.3782\\
Variant & Others  & 13 & 0.2373 & 0.2993 & 0.4281 & 0.4442\\
\midrule
Variant & All     & 36 & 0.4325 & 0.4451 & 0.3127 & 0.3876\\
\bottomrule
\end{tabular}
\par\vspace{3pt}
\parbox{\linewidth}{\footnotesize
LOTO removes only the test category's references. Scores are not
percentages of shared constraints or accuracy measures; no universal
high/low threshold is assumed.}
\end{table}

\medskip
\textbf{Results and interpretation.}
The 101-instance same-category median and mean are $0.5065$ and
$0.4632$, while the LOTO median and mean are $0.4424$ and $0.4745$.
For the 36 variants these four values are $0.4325$, $0.4451$, $0.3127$,
and $0.3876$.
The category values vary substantially, so category membership does not
make formulations interchangeable. AP's low same-category score is
partly explained by continuous test variables versus a binary AP
reference. NRM is also affected by a synthetic shared-dispatch row in
the revised reference; removing that row from a diagnostic copy restores
the 25 NRM scores to $1$. These observations qualify the interpretation
of the graph scores and do not by themselves prove new business
requirements in the test instances.

Withholding only the designated category's references gives Mixture
median and mean LOTO scores of $0.4057$ and $0.4327$ in the original set,
and $0.2520$ and $0.3782$ for the variants. The corresponding Others
values are $0.3615$, $0.4070$ and $0.4281$, $0.4442$. This tests reliance
on same-category Ref-Data while keeping Mixture and Others separate.

The WL representation omits objective coefficients, matrix values and
signs, right-hand sides, and bound values. A high score therefore does
not establish model equivalence, and a low score does not establish
business-level novelty or robustness. The structural LOTO result is a
similarity diagnostic, not a solver-accuracy experiment.

\medskip
\textbf{Complementary performance evidence.}
Separate saved evaluations test behavior under changed tasks and
workflow support. The 36-variant evaluation matched the reference
objective on 29/36 instances ($80.56\%$). In the saved GPT-4.1
examples-and-route withholding experiment, 88/101 instances matched
($87.13\%$), versus $94.06\%$ for that experiment's saved full baseline.
In forced routing, 431/505 non-designated route assignments matched
($85.35\%$); these are repeated assignments to the 101 problems. All
failures remain in the denominators, and objective matching is only a
performance indicator. These records come from their respective saved
implementation snapshots and are not pooled as one common-version
comparison.

Together, the structural and performance evidence supports resilience
beyond dependence on one intact same-category example under the tested
interventions. It does not establish universal generalization to
arbitrary new constraints, objectives, or operational requirements.

\medskip
\textbf{Changes in the manuscript.}
We added a new subsection, ``Formulation-Level Structural Comparisons
between \refdata{} and \testdata{}'', which briefly describes how the
structural similarity is computed and reports the final results.
\endgroup

\begingroup\color{blue}
\subsection{\LK{Formulation-Level Structural Comparisons between
\refdata{} and \testdata{}}}
To complement the textual comparison, we represent each LP as a typed
variable--constraint bipartite graph and compare the graphs with a
two-round Weisfeiler--Lehman (WL) subtree feature map
\citep{Gasse2019,Shervashidze2011}. Singleton rows are absorbed into
effective variable bounds, which remain active in the optimization model
but are excluded from the graph; variable domains and all remaining rows
are retained. The graph records variable--constraint incidence and row or
variable types, while omitting names, objective terms, coefficients,
right-hand sides, and bound values. We concatenate label counts from
rounds $0$--$2$ and use cosine similarity. For each test formulation, we
report the maximum similarity to its same-category references
($S_{\rm same}$) and after removing all references with that category
($S_{\rm LOTO}$). The table reports medians and means of these per-test
maxima; LOTO removes only the test category's Ref-Data.

\begin{table}[H]
\centering\small
\caption{Formulation-level WL similarity between \refdata{} and the
benchmarks. Effective bounds are retained as variable restrictions and
excluded from the graph; all scores use cumulative rounds $0$--$2$.}
\label{tab:manuscript-wl-summary}
\setlength{\tabcolsep}{6pt}
\begin{tabular}{@{}lrrrrr@{}}
\toprule
Benchmark set & Tests & Median $S_{\rm same}$ & Mean $S_{\rm same}$ &
Median $S_{\rm LOTO}$ & Mean $S_{\rm LOTO}$\\
\midrule
\testdata{} & 101 & 0.5065 & 0.4632 & 0.4424 & 0.4745\\
Task variants & 36 & 0.4325 & 0.4451 & 0.3127 & 0.3876\\
\bottomrule
\end{tabular}
\end{table}

These results show that the available Ref-Data are not interchangeable
with the benchmark formulations under the selected structural encoding;
the category-level results vary substantially, and the LOTO scores remain
structural diagnostics rather than accuracy or universal generalization
measures.
\endgroup
